\documentclass[10pt,a4paper]{article}

% Standard packages.
\usepackage{lmodern}
\usepackage[T1]{fontenc}
\usepackage[utf8]{inputenc}
\usepackage[margin=1in]{geometry}
\usepackage{booktabs}
\usepackage{longtable}
\usepackage{array}
\usepackage{listings}
\usepackage{xcolor}
\usepackage[most]{tcolorbox}
\usepackage{microtype}
\usepackage[hidelinks]{hyperref}
\usepackage{parskip}

% Define some colours, got from eye dropper on the user guide for SLF.
\definecolor{shadecolor}{RGB}{248,248,248}
\definecolor{codekey}{RGB}{5,41,125}
\definecolor{codestring}{RGB}{64,112,160}
\definecolor{codecomment}{RGB}{97,161,176}
\definecolor{codebool}{RGB}{0,112,33}

% Stop -- turning into en-dash in plain text, for CLI commands.
\DisableLigatures[-]{encoding = T1, family = tt*}

% Allows underscores to override line breaks in longtable.
\let\origunderscore\_
\renewcommand{\_}{\origunderscore\penalty4000\relax}
\newcolumntype{L}[1]{>{\raggedright\arraybackslash}p{#1}}

% Loosen tolerance for line breaking.
\tolerance=1200
\setlength{\emergencystretch}{3em}

% Settings to make code blocks look like Python.
\lstdefinestyle{py}{
    basicstyle=\ttfamily\small,
    commentstyle=\color{codecomment}\itshape,
    morecomment=[l]{\#},
    morestring=[b]",
    morestring=[b]',
    stringstyle=\color{codestring},
    alsoletter={_},
    emph={Stratos, normal, oblique, shock_angle,
    interpolate_to_length, gasFile, chemFile, kinFile,
    x_end, cells, u1, massf1, T1, p1, rho1, theta, beta,
    verbose, end_type, max_iter, tol, no_cells,
    target_length, columns, reference},
    emphstyle={\color{codekey}},
    emph=[2]{True,False},
    emphstyle=[2]{\color{codebool}\bfseries},
    frame=none,
    showstringspaces=false,
    columns=fullflexible,
    keepspaces=true,
    breaklines=true,
}
\lstdefinestyle{plain}{
    basicstyle=\ttfamily\small,
    emph={},
    emph=[2]{},
    frame=none,
    columns=fullflexible,
    keepspaces=true,
    breaklines=true,
}
\lstset{style=py}

% Put a shaded box around listings.
\tcolorboxenvironment{lstlisting}{
    breakable,
    colback=shadecolor,
    colframe=shadecolor,
    boxrule=0pt,
    arc=4pt, outer arc=4pt,
    boxsep=0pt,
    left=6pt, right=6pt, top=6pt, bottom=6pt,
    grow to left by=6pt, grow to right by=6pt,
    before skip=8pt, after skip=8pt,
}

%+++++++++++++++++++++++++++++++++++++++++++++
%             Preamble ends here
%+++++++++++++++++++++++++++++++++++++++++++++

\title{STRATOS: Reference Manual}
\author{\textit{Taine J. Rossini}}
\date{Version 1.0.0.}

\begin{document}
\maketitle

\section*{Abstract}

This document lists every argument recognised by \texttt{stratos}, the post-shock relaxation solver, and the returned data.

\section*{Usage}

Stratos has no command line interface. It is a single class, imported into a job script.

\begin{lstlisting}
from stratos import Stratos

st = Stratos('air-5sp.gas', 'air-5sp.chem', '', 0.1, 1000, 6075.1,
             {"N2": 0.767, "O2": 0.233}, 365.5, rho1=0.02322)

columns = st.normal()
\end{lstlisting}

A run has three parts.

\begin{itemize}
    \item Construct a \texttt{Stratos} object from the gas-models, domain, and pre-shock freestream.
    \item Call one solution method, \texttt{normal()}, \texttt{oblique()} or \texttt{shock\_angle()}, to march along the streamline.
    \item Use the returned dictionary. Nothing is written to disk and nothing is plotted; every method hands back a dictionary of equal-length numpy arrays for the script to do with as it likes.
\end{itemize}

One object can be used for as many runs as needed. The freestream is fixed at construction, so a sweep over shock angles is a loop over the method, not over the constructor.

\section*{Constructor arguments}

The arguments are positional, in this order, and none of the first eight has a default.

\begin{longtable}{@{}lL{22mm}L{102mm}@{}}
    \toprule
    \textbf{Argument} & \textbf{Type} & \textbf{Description} \\
    \midrule
    \endfirsthead
    \toprule
    \textbf{Argument} & \textbf{Type} & \textbf{Description} \\
    \midrule
    \endhead
    \bottomrule
    \endfoot
    \endlastfoot
    \texttt{gasFile} & string & Path to the gas-model file. \\
    \addlinespace
    \texttt{chemFile} & string & Path to the reaction file. \\
    \addlinespace
    \texttt{kinFile} & string & Path to the energy-exchange file, for two-temperature runs. Pass \texttt{""} on a one-temperature gas-model. \\
    \addlinespace
    \texttt{x\_end} & float & Length of the marching domain, m. \\
    \addlinespace
    \texttt{cells} & int & Number of points the domain is divided into. \\
    \addlinespace
    \texttt{u1} & float & Freestream velocity, m/s. \\
    \addlinespace
    \texttt{massf1} & dict or list & Freestream mass fractions. A dictionary is keyed by species name. \\
    \addlinespace
    \texttt{T1} & float & Freestream temperature, K. \\
    \bottomrule
\end{longtable}

In addition, exactly one of \texttt{p1} or \texttt{rho1} must be supplied as a keyword.

\begin{center}
	\begin{tabular}{@{}lllL{88mm}@{}}
	\toprule
	\textbf{Argument} & \textbf{Type} & \textbf{Default} & \textbf{Description} \\
	\midrule
	\texttt{p1} & float & \texttt{None} & Freestream pressure, Pa. \\
	\addlinespace
	\texttt{rho1} & float & \texttt{None} & Freestream density, kg/m$^3$. \\
	\bottomrule
	\end{tabular}
\end{center}

If both are given, \texttt{p1} is used. If neither is, the construction fails.

\section*{Solution methods}

\subsection*{\texttt{normal()}}

Applies the normal-shock jump conditions to the freestream and marches from the immediate post-shock state to \texttt{x\_end}. Returns the column dictionary.

\begin{lstlisting}
columns = st.normal()
\end{lstlisting}

\subsection*{\texttt{oblique(theta)}}

Solves the shock over a wedge of half-angle \texttt{theta}. The shock angle from the frozen relation is only a starting guess; a secant iteration adjusts it until the flow angle at the end of the domain matches the wedge. Returns a tuple of the column dictionary and the converged shock angle, in radians.

\begin{longtable}{@{}lL{16mm}L{22mm}L{80mm}@{}}
	\toprule
	\textbf{Argument} & \textbf{Type} & \textbf{Default} & \textbf{Description} \\
	\midrule
	\endfirsthead
	\toprule
	\textbf{Argument} & \textbf{Type} & \textbf{Default} & \textbf{Description} \\
	\midrule
	\endhead
	\bottomrule
	\endfoot
	\endlastfoot
	\texttt{theta} & float & required & Wedge half-angle, radians. \\
	\addlinespace
	\texttt{verbose} & bool & \texttt{True} & Prints the residual, last flow angle and shock angle each iteration. \\
	\addlinespace
	\texttt{end\_type} & string & \texttt{'x\_arc'} & Which column \texttt{x\_end} is measured along, and which the result is interpolated onto. \\
	\addlinespace
	\texttt{max\_iter} & int & \texttt{30} & Maximum secant iterations for the shock angle iteration. \\
	\addlinespace
	\texttt{tol} & float & \texttt{1e-5} & Convergence tolerance on the flow angle, radians. \\
	\bottomrule
\end{longtable}

\begin{lstlisting}
columns, beta = st.oblique(np.deg2rad(15.0), verbose=False)
\end{lstlisting}

\subsection*{\texttt{shock\_angle(beta)}}

Marches the streamline behind a shock of known angle \texttt{beta}, in radians. This is the method for tracing a shock whose shape is already known. Returns the column dictionary.

\begin{lstlisting}
columns = st.shock_angle(beta)
\end{lstlisting}

The tangential velocity is preserved across the shock and the normal component marched, so this and \texttt{oblique()} share their post-processing and return the same columns, less \texttt{x\_p}.

\subsection*{\texttt{interpolate\_to\_length(no\_cells, target\_length, columns, reference)}}

Interpolates a result onto an evenly spaced grid. It can be used to re-grid a result afterwards, for example onto residence time rather than distance.

\begin{center}
	\begin{tabular}{@{}lL{18mm}L{100mm}@{}}
	\toprule
	\textbf{Argument} & \textbf{Type} & \textbf{Description} \\
	\midrule
	\texttt{no\_cells} & int & Number of points in the new grid. \\
	\addlinespace
	\texttt{target\_length} & float & Upper limit of the new grid, in the units of \texttt{reference}. \\
	\addlinespace
	\texttt{columns} & dict & The dictionary to re-grid. Every column in it is interpolated. \\
	\addlinespace
	\texttt{reference} & string & The column that the new grid spans. \\
	\bottomrule
	\end{tabular}
\end{center}

\begin{lstlisting}
st.interpolate_to_length(1000, 2e-7, columns, 't')
\end{lstlisting}

\section*{Returned data}
\label{sec:returned}

Every method returns a dictionary of numpy arrays, all \texttt{cells} long.

\subsection*{From \texttt{normal()}}

\begin{longtable}{@{}lL{22mm}L{102mm}@{}}
	\toprule
	\textbf{Key} & \textbf{Units} & \textbf{Description} \\
	\midrule
	\endfirsthead
	\toprule
	\textbf{Key} & \textbf{Units} & \textbf{Description} \\
	\midrule
	\endhead
	\bottomrule
	\endfoot
	\endlastfoot
	\texttt{x} & m & Distance behind the shock. \\
	\addlinespace
	\texttt{t} & s & Time behind the shock, integrated along the streamline. \\
	\addlinespace
	\texttt{u} & m/s & Velocity. \\
	\addlinespace
	\texttt{T\_tr} & K & Translational-rotational temperature. \\
	\addlinespace
	\texttt{T\_ve} & K & Vibro-electronic temperature. Two-temperature runs only. \\
	\addlinespace
	\texttt{p} & Pa & Pressure. \\
	\addlinespace
	\texttt{rho} & kg/m$^3$ & Density. \\
	\addlinespace
	\texttt{e\_tr} & J/kg & Translational-rotational internal energy. \\
	\addlinespace
	\texttt{e\_ve} & J/kg & Vibro-electronic internal energy. Two-temperature runs only. \\
	\addlinespace
	\texttt{gamma} & - & Ratio of specific heats. \\
	\addlinespace
	\texttt{R} & J/kg/K & Gas constant of the mixture. \\
	\addlinespace
	\texttt{Y\_<species>} & - & Mass fraction, one column per species. \\
	\addlinespace
	\texttt{X\_<species>} & - & Mole fraction, one column per species. \\
	\bottomrule
\end{longtable}

On a one-temperature gas-model the two vibro-electronic columns are absent.

\subsection*{From \texttt{oblique()} and \texttt{shock\_angle()}}

Both carry every column above, except that \texttt{x} is removed in favour of the shock-aligned and Cartesian coordinates below, and \texttt{u} is redefined as the velocity magnitude.

\begin{longtable}{@{}lL{22mm}L{102mm}@{}}
	\toprule
	\textbf{Key} & \textbf{Units} & \textbf{Description} \\
	\midrule
	\endfirsthead
	\toprule
	\textbf{Key} & \textbf{Units} & \textbf{Description} \\
	\midrule
	\endhead
	\bottomrule
	\endfoot
	\endlastfoot
	\texttt{x\_n} & m & Distance normal to the shock. \\
	\addlinespace
	\texttt{x\_t} & m & Distance tangential to the shock. \\
	\addlinespace
	\texttt{x\_arc} & m & Distance along the streamline, the default \texttt{end\_type}. \\
	\addlinespace
	\texttt{x\_coord} & m & Cartesian $x$, aligned with the freestream, with the origin where the streamline crosses the shock. \\
	\addlinespace
	\texttt{y\_coord} & m & Cartesian $y$, normal to the freestream. \\
	\addlinespace
	\texttt{x\_p} & m & Distance along the wedge surface. \texttt{oblique()} only. \\
	\addlinespace
	\texttt{u\_n} & m/s & Velocity normal to the shock, which is what is marched. \\
	\addlinespace
	\texttt{u\_t} & m/s & Velocity tangential to the shock, constant across it. \\
	\addlinespace
	\texttt{u} & m/s & Velocity magnitude, not the normal component. \\
	\addlinespace
	\texttt{flow\_angle} & rad & Angle of the streamline, measured from the freestream. \\
	\bottomrule
\end{longtable}

\section*{A complete job script}

A one-temperature normal shock, marched a metre downstream.

\begin{lstlisting}
# Load libraries.
from stratos import Stratos
import numpy as np

# Construct the solver.
st = Stratos('air-7sp.gas', 'air-7sp.chem', '', 1.0, 100000, 4273.64,
             {"N2": 0.78, "O2": 0.22}, 300, 133.3)

columns = st.normal()

# Save the result to a CSV file.
np.savetxt('profile.csv', np.column_stack(list(columns.values())),
           delimiter=',', header=','.join(columns.keys()), comments='')
\end{lstlisting}

A sweep of wedge angles on the same object.

\begin{lstlisting}
# Load libraries.
from stratos import Stratos
import numpy as np

# Construct the solver.
st = Stratos('air-5sp.gas', 'air-5sp.chem', 'air-5sp.kin', 0.1, 1000, 6075.1,
             {"N2": 0.767, "O2": 0.233}, 365.50, rho1=0.02322)

results = {}

for angle in [5, 10, 15, 20, 25, 30, 35, 40]:
    results[angle], _ = st.oblique(np.deg2rad(angle), verbose=False)

# Print the results at the end of the domain.
for angle, columns in results.items():
    print(angle, columns['T_tr'][-1])
\end{lstlisting}

\end{document}
